周报里有这样一句结论：``加上新特征后，模型在验证集上稳定提升。''一位挑剔的读者会连问三句：稳定指重复了多少次？换一批验证集还成立吗？提升会不会来自同期的别的改动？三问之后，原来那句话要么被补成一句带条件的断言，要么当场垮掉。

证明做的是同一件事，只是把标准提到顶格。它要求结论从已经承认的东西里\emph{不得不}成立，并且这条链子上的每一环，都经得起一个完全不信任你的人逐环追问。

有个流传很广的误解：会不会证明取决于见过多少花招。真正起作用的判断力朴素得多——先看清结论的逻辑形状，形状规定了什么样的证据才算数；剩下的工作，是在几条标准路线里挑一条最短的走。本章就按这个顺序展开。

本章的例子都刻意选得短。要练的不是代数熟练度，而是两句追问：这个结论要求我交付什么？眼前这一步凭什么合法？

\subsection{怎么读这一章}

四篇构成一条主线，建议第一次按顺序读：

\begin{enumerate}
\tightlist
\item \hyperref[sec:01-Mathematical-Language/03-Proof-Techniques/01]{``从命题形状开始证明''}：把``该怎么证''翻译成``该交付什么证据''，并看清反例必须长什么样；
\item \hyperref[sec:01-Mathematical-Language/03-Proof-Techniques/02]{``逆否、反证与分类讨论''}：正面走不通时换入口，以及每种换法各自需要什么条件才启动得起来；
\item \hyperref[sec:01-Mathematical-Language/03-Proof-Techniques/03]{``存在性与唯一性''}：把一个符号拆成两笔独立的账，并区分``能算出来''和``只知道有''；
\item \hyperref[sec:01-Mathematical-Language/03-Proof-Techniques/04]{``怎样发现并修复错误证明''}：错误证明的几种固定形态，以及它们在数据工作里的等价物。
\end{enumerate}

如果眼下的困惑是``我写的东西到底算不算证明''，可以先跳到第四篇，用它反过来照前三篇。第四篇本身也值得单独重读——它训练的是审阅别人结论的能力，用途远超数学作业。

\subsection{读完后应发生什么变化}

目标不是背下手法清单，而是拿到任意一句待证断言时，能立刻说出三件事：它的逻辑形状是什么，因而合格的证据长什么样；如果它是假的，反例必须满足哪些条件；证明里最可能藏错的是哪一步。这三个动作在后面每一章都会被反复调用，尤其是在读定理条件和判断实验结论是否可信的时候。
